\documentclass[11pt]{article}

% ─────────────────────────────────────────────────────────────────────────────
% Packages
% ─────────────────────────────────────────────────────────────────────────────
\usepackage[T1]{fontenc}
\usepackage[utf8]{inputenc}
\usepackage{lmodern}
\usepackage[margin=1in]{geometry}
\usepackage{amsmath,amssymb,amsthm,mathtools}
\usepackage{bm}
\usepackage{mathrsfs}
\usepackage{enumitem}
\usepackage{array}
\usepackage{booktabs}
\usepackage{longtable}
\usepackage{tabularx}
\usepackage{hyperref}
\usepackage[nameinlink,capitalise]{cleveref}
\usepackage{xcolor}
\usepackage{listings}
\usepackage{setspace}
\usepackage{titlesec}

\hypersetup{
  colorlinks=true,
  linkcolor=blue!60!black,
  citecolor=blue!60!black,
  urlcolor=blue!60!black,
  pdftitle={Multiplicity Open Core and the Core Multiplicity Operator},
  pdfauthor={},
  pdfsubject={Multiplicity Theory},
  pdfkeywords={Multiplicity Theory, operator theory, prime-graded Hilbert bundle, open core, lawfulness projector}
}

\onehalfspacing

% ─────────────────────────────────────────────────────────────────────────────
\theoremstyle{plain}
\newtheorem{theorem}{Theorem}[section]
\newtheorem{lemma}[theorem]{Lemma}
\newtheorem{proposition}[theorem]{Proposition}
\newtheorem{corollary}[theorem]{Corollary}
\newtheorem{conjecture}[theorem]{Conjecture}
\newtheorem{claim}[theorem]{Claim}

\theoremstyle{definition}
\newtheorem{definition}[theorem]{Definition}
\newtheorem{example}[theorem]{Example}
\newtheorem{construction}[theorem]{Construction}
\newtheorem{assumption}[theorem]{Assumption}
\newtheorem{notation}[theorem]{Notation}

\theoremstyle{remark}
\newtheorem{remark}[theorem]{Remark}
\newtheorem{note}[theorem]{Note}
\newtheorem{observation}[theorem]{Observation}
\newtheorem{warning}[theorem]{Warning}

─────────────────────────────────────────────────────────────────────────────
% PROOF CUSTOMIZATION
% ─────────────────────────────────────────────────────────────────────────────
% Named proofs (e.g. "Proof of Theorem 3.1")
\newenvironment{namedproof}[1]{%
  \begin{proof}[Proof of #1]}{\end{proof}}

% Proof sketch
\newenvironment{proofsketch}{%
  \begin{proof}[Proof sketch]}{\end{proof}}

% QED symbol — solid black square (AMS standard)
\renewcommand{\qedsymbol}{\ensuremath{\blacksquare}}
% Macros
% ─────────────────────────────────────────────────────────────────────────────
\newcommand{\Hbundle}{\mathcal{H}}
\newcommand{\Pset}{\mathbb{P}}
\newcommand{\R}{\mathbb{R}}
\newcommand{\C}{\mathbb{C}}
\newcommand{\Om}{\boldsymbol{\Omega}}
\newcommand{\PCSL}{P_{\mathrm{CSL}}}
\newcommand{\PFix}{P_{\mathrm{Fix}(\Xi)}}
\newcommand{\norm}[1]{\left\lVert #1 \right\rVert}
\newcommand{\abs}[1]{\left\lvert #1 \right\rvert}
\newcommand{\ip}[2]{\left\langle #1, #2 \right\rangle}
\newcommand{\sot}{\mathrm{s\text{-}lim}}
\newcommand{\HS}{\mathrm{HS}}
\newcommand{\id}{\mathrm{id}}

% ─────────────────────────────────────────────────────────────────────────────
% Listings
% ─────────────────────────────────────────────────────────────────────────────
\definecolor{codebg}{RGB}{248,248,248}
\definecolor{codekw}{RGB}{0,70,140}
\definecolor{codestr}{RGB}{163,21,21}
\definecolor{codecom}{RGB}{0,128,0}

\lstdefinestyle{pythonstyle}{
  language=Python,
  basicstyle=\ttfamily\small,
  backgroundcolor=\color{codebg},
  keywordstyle=\color{codekw}\bfseries,
  stringstyle=\color{codestr},
  commentstyle=\color{codecom},
  showstringspaces=false,
  frame=single,
  breaklines=true,
  tabsize=4
}

% ─────────────────────────────────────────────────────────────────────────────
% Title
% ─────────────────────────────────────────────────────────────────────────────
\title{\textbf{Multiplicity Open Core and the Core Multiplicity Operator}\\
\large A Unification Report on the Developments of the Present Thread}
\author{Citizen Gardens \\ The Foundation of Multiplicity}
\date{April 2026}

\begin{document}
\maketitle

\begin{abstract}
This report consolidates the main conceptual, mathematical, and architectural developments produced in the present thread. The central conclusion is that the relevant sectoral engines---including Four-Prime Newton Holography, CRMF, the Genomic Analogy of an AGI Immune System, and future sectors such as Moonshine, CEQG Track C, UME, and the Langlands Prism---should not be collapsed into a premature claim of literal identity, but instead organized through a disciplined abstraction: the \emph{Multiplicity Open Core}. This open core isolates the invariant simulator pattern shared by lawful multiplicity dynamics: visible indexed generators, a latent stabilizer, and a lawful projection, with an optional certification layer.

A second major development is the introduction of the \emph{Core Multiplicity Operator} \(\Om\), defined on a prime-graded Hilbert bundle and intended as a canonical mother operator from which subsystem-specific operators arise by lawful restriction. This report formalizes the open core, defines the operator-theoretic setting of \(\Om\), states the main specialization theorem, presents explicit operator norm bounds, gives a Banach fixed-point proof for the \(\Lambda_m\)-sector, and supplies implementation-ready Python skeletons. The resulting framework provides a mathematically coherent path from philosophical unification to modular software architecture and eventual formal verification.
\end{abstract}
\newpage
\tableofcontents
\newpage

% ─────────────────────────────────────────────────────────────────────────────
\section{Executive Summary}
% ─────────────────────────────────────────────────────────────────────────────

The thread converged on two mutually compatible abstractions.

First, the three principal engines under discussion remain \emph{sectorally autonomous}: Four-Prime Newton Holography, CRMF, and the Genomic Analogy of an AGI Immune System each preserve their own proofs, observables, and engineering or clinical claims. Their unity is not enforced by collapsing them into one rhetorical object, but by extracting a reusable simulator substrate: the \emph{Multiplicity Open Core}.

Second, the software architecture can be made precise through a canonical operator layer, the \emph{Core Multiplicity Operator} \(\Om\), acting on a prime-graded Hilbert bundle. Sectoral engines become lawful restrictions or specializations of \(\Om\), rather than disconnected implementations.

The strongest architectural recommendations produced in the thread are:

\begin{enumerate}[label=\textbf{R\arabic*.}]
    \item Preserve sectoral autonomy while elevating common structure into a formal open core.
    \item Implement \texttt{omega\_operator.py} as a thin synthesis layer rather than a rewrite of existing modules.
    \item Distinguish operator construction, admissibility, specialization, and certification as separate interfaces.
    \item Treat lawfulness as a first-class invariant, not as an informal design aspiration.
    \item Use the frame-flat specialization as the first validation target before advancing to modular or stochastic sectors.
\end{enumerate}

% ─────────────────────────────────────────────────────────────────────────────
\section{Background and Conceptual Development}
% ─────────────────────────────────────────────────────────────────────────────

\begin{definition}[Multiplicity Theory]
Multiplicity Theory is a prime-indexed, recursively stable mathematical framework in which multiplicity---understood as recurrence structured through prime decomposition---serves as the basis for modeling nonlinear and emergent phenomena across mathematical, physical, computational, and cognitive domains.
\end{definition}

The thread developed this theory in two complementary directions. One direction was conceptual: how to preserve the integrity of independent engines while identifying their invariant structural pattern. The other was technical: how to implement that invariant pattern in code without forcing rewrites of existing modules.

The resulting invariant pattern is:

\begin{center}
\textbf{Visible Indexed Generators} \(+\) \textbf{Latent Stabilizer} \(+\) \textbf{Lawful Projection}.
\end{center}

This pattern is general enough to cover:
\begin{itemize}
    \item sectoral physical readouts,
    \item resonance-governed biological or biosemantic systems,
    \item self/non-self discrimination architectures,
    \item operator restrictions in modular, stochastic, and functorial settings.
\end{itemize}

% ─────────────────────────────────────────────────────────────────────────────
\section{Multiplicity Open Core Specification v0.1}
% ─────────────────────────────────────────────────────────────────────────────

\begin{definition}[Multiplicity Open Core]
The \emph{Multiplicity Open Core} is the minimal simulator pattern sufficient to generate lawful multiplicity dynamics. It consists of three required slots:
\begin{enumerate}
    \item \textbf{Generators},
    \item \textbf{Stabilizer},
    \item \textbf{Projection},
\end{enumerate}
together with an optional \textbf{Certification Layer}.
\end{definition}

\subsection{Generators}

The generators encode the visible indexed dynamics of a sector.

\begin{equation}
\Xi(t) = \sum_{p \in P_N} a_p(t)\, U_p(t).
\label{eq:generators}
\end{equation}

Here \(P_N\) is a finite prime index set, \(a_p(t)\) are sector-dependent amplitudes or coefficients, and \(U_p(t)\) are prime-indexed operators.

\paragraph{Interpretation.}
The generators describe the \emph{apparent} or observable dynamic structure. They do not, by themselves, guarantee lawful recursion or stable emergence.

\subsection{Stabilizer}

The stabilizer imposes a hidden constraint that enforces boundedness, redundancy, fixed-point admissibility, or holographic consistency.

\begin{equation}
[\Om, P] \approx 0,
\label{eq:stabilizer-commutator}
\end{equation}

or, more generally, through a resonance functional \(R_t\) and contraction certificate \(\gamma_t\).

\paragraph{Interpretation.}
The stabilizer is the latent lawfulness layer. It is what turns a family of visible dynamics into a controlled multiplicity process.

\subsection{Projection}

The projection maps the universal core onto an observable domain-specific readout.

\begin{equation}
R_s(\Om) \longrightarrow \Om_s,
\label{eq:projection}
\end{equation}

where \(R_s\) is a lawful sector restriction and \(\Om_s\) is the specialized operator.

\paragraph{Interpretation.}
Projection is not arbitrary output formatting. It is a lawful restriction or readout map that preserves enough structure for the specialization to remain meaningful.

\subsection{Certification Layer}

The certification layer is optional at the minimal specification level, but necessary for machine-checkable engines.

Typical certification conditions include:
\begin{align}
\norm{[\Om_s, P]} &\le \varepsilon, \label{eq:c1}\\
\gamma_t &\le 1, \label{eq:c2}\\
R_t &\in [\rho_{\min},\rho_{\max}], \label{eq:c3}\\
\mathcal{M}_t &= H(\mathcal{M}_{t-1}\Vert \Om_{s,t}). \label{eq:c4}
\end{align}

\subsection{Sectoral Mapping}

\begin{longtable}{>{\raggedright\arraybackslash}p{0.18\textwidth} >{\raggedright\arraybackslash}p{0.24\textwidth} >{\raggedright\arraybackslash}p{0.24\textwidth} >{\raggedright\arraybackslash}p{0.24\textwidth}}
\toprule
\textbf{Sector} & \textbf{Generators} & \textbf{Stabilizer} & \textbf{Projection} \\
\midrule
Four-Prime Newton Holography &
Three visible primes encoding apparent Newtonian structure &
Fourth latent prime as holographic stabilizer &
Delay embeddings and prime-Cantor fractal readout \\
CRMF &
Prime-indexed pathway operators or adjacency operators &
Resonance functional, CSC clamp, contraction condition &
Sparse PMDM and FWHT codon-contrast field \\
AGI Immune Analogy &
ADRs or governance loci as indexed units &
L0/L1/L2 hierarchy and self-audit layer &
Witness-gate and self/non-self predicate \\
Meta-Relativity &
Universal spectral blocks \(A,B,E\) &
Lawfulness projector in curvature-free sector &
Frame-flat restriction \\
Moonshine &
Modular or Hecke-indexed iteration operators &
Fixed-point modular admissibility &
Moonshine restriction on upper half-plane \\
CEQG Track C &
Noise and dissipation operators &
Stochastic admissibility and bounded dissipation &
Noise-kernel specialization \\
UME &
Neuromorphic or ACE-projected generators &
ACE admissibility &
Cognitive projection \\
Langlands Prism &
Automorphic or functorial blocks &
Functorial compatibility &
Prismatic control-layer image \\
\bottomrule
\end{longtable}

% ─────────────────────────────────────────────────────────────────────────────
\section{Core Multiplicity Operator}
% ─────────────────────────────────────────────────────────────────────────────

\subsection{Prime-Graded Hilbert Bundle}

The canonical state space is the prime-graded Hilbert bundle
\begin{equation}
\Hbundle
=
\bigoplus_{p \in \Pset}
\mathcal{H}_p \otimes L^2(\R) \otimes \C^d.
\label{eq:hilbert-bundle}
\end{equation}

This expression encodes three layers:
\begin{itemize}
    \item a prime-indexed grading \(\mathcal{H}_p\),
    \item a time or signal factor \(L^2(\R)\),
    \item a finite internal sector \(\C^d\).
\end{itemize}

\subsection{Operator Decomposition}

The core operator is defined by
\begin{equation}
\Om = D_\sigma + K_{\alpha,h} + C + \Xi.
\label{eq:omega-decomposition}
\end{equation}

The four blocks are interpreted as follows:
\begin{itemize}
    \item \(D_\sigma\): diagonal prime-weighted block,
    \item \(K_{\alpha,h}\): windowed off-diagonal or coupling block,
    \item \(C\): time-sieve multiplier,
    \item \(\Xi\): internal or recursive block.
\end{itemize}

\subsection{Lawfulness Projector}

Let \(\Pi_p\) denote the orthogonal projection onto the \(p\)-th prime fibre. Define
\begin{equation}
P_{\Pi}
=
\sot_{N \to \infty}\;
\sum_{p \le p_N} \Pi_p,
\label{eq:Ppi}
\end{equation}
and let \(P_{\mathrm{Fix}(\Xi)}\) be the fixed-point spectral projector of \(\Xi\). Then the lawfulness projector is
\begin{equation}
\PCSL = P_{\Pi} \wedge P_{\mathrm{Fix}(\Xi)}.
\label{eq:pcsl}
\end{equation}

\begin{remark}
The operator \(\PCSL\) is intended to encode lawful admissibility across the open core. In implementation terms it is built once and reused in every specialization pathway.
\end{remark}

% ─────────────────────────────────────────────────────────────────────────────
\section{Main Mathematical Statements}
% ─────────────────────────────────────────────────────────────────────────────

\begin{theorem}[Commutation Invariant]
Assume the window function satisfies \(\widehat{h} \in L^1(\R)\), the block \(C\) acts only on the \(L^2(\R)\)-factor, and \(P_{\mathrm{Fix}(\Xi)}\) is the spectral projection onto the fixed-point subspace of \(\Xi\). Then
\begin{equation}
[\Om,\PCSL]=0.
\label{eq:commutation-invariant}
\end{equation}
\end{theorem}

\begin{proof}[Proof sketch]
The diagonal block \(D_\sigma\) commutes with every fibre projection \(\Pi_p\) by construction. Under the stated window condition, the coupling block \(K_{\alpha,h}\) is bounded and fibre-compatible. The block \(C\) acts on a distinct tensor factor and therefore commutes with prime-fibre projections. Finally, \(P_{\mathrm{Fix}(\Xi)}\) is a spectral projection of \(\Xi\), so it commutes with \(\Xi\). The meet of commuting projections remains compatible with any operator commuting with each factor, yielding \eqref{eq:commutation-invariant}.
\end{proof}

\begin{theorem}[Sector Specialization]
Let \(S\) be the set of admissible sectors and let \(R_s\) denote a lawful restriction map for \(s \in S\). Then for each admissible sector \(s\), there exists a specialized operator
\[
\Om_s = R_s(\Om)
\]
such that:
\begin{enumerate}[label=(\roman*)]
    \item \(\Om_s\) acts on a \( \PCSL \)-invariant subspace,
    \item \(\norm{\Om_s} \le \norm{\Om}\),
    \item sector composition is functorial whenever the corresponding sector meet exists.
\end{enumerate}
\end{theorem}

\begin{proof}[Proof sketch]
The restriction \(R_s\) is implemented as compression to a sector-invariant subbundle of \(\Hbundle\). Compression to an invariant subspace cannot increase the operator norm. Functoriality follows from compatibility of inclusions and orthogonal projections between nested sector spaces.
\end{proof}

\begin{theorem}[Global Norm Bound]
Suppose \(C = \Lambda_m \circ M_\varphi\) with \(m \in L^\infty(\R)\) and \(\varphi \in L^\infty(\R)\), and assume \(\Xi\) is self-adjoint with spectral radius \(r(\Xi)\). Then
\begin{equation}
\norm{\Om}
\le
|\sigma| + \norm{\widehat{h}}_{L^1}
+ \norm{m}_{L^\infty}\norm{\varphi}_{L^\infty}
+ r(\Xi).
\label{eq:global-bound}
\end{equation}
\end{theorem}

\begin{proof}[Proof sketch]
Apply the triangle inequality to \eqref{eq:omega-decomposition}. The norm of \(D_\sigma\) is \(|\sigma|\). The bound for \(K_{\alpha,h}\) is controlled by \(\norm{\widehat{h}}_{L^1}\). The block \(C\) obeys the operator-norm product bound, and the norm of self-adjoint \(\Xi\) equals its spectral radius.
\end{proof}

\begin{theorem}[Banach Fixed-Point Certificate for \(\Lambda_m\)]
Let \((\mathcal{X},\norm{\cdot})\) be a Banach space and suppose the recursive update map \(\Phi:\mathcal{X}\to\mathcal{X}\) satisfies
\begin{equation}
\norm{\Phi(\psi)-\Phi(\phi)} \le \gamma \norm{\psi-\phi}
\quad\text{for all } \psi,\phi\in\mathcal{X},
\label{eq:banach-contract}
\end{equation}
with \(0\le \gamma <1\). Then \(\Phi\) has a unique fixed point and the iterates converge geometrically:
\begin{equation}
\norm{\psi_n-\psi^\star} \le \gamma^n \norm{\psi_0-\psi^\star}.
\label{eq:banach-geometric}
\end{equation}
\end{theorem}

\begin{proof}
This is the Banach fixed-point theorem applied to the \(\Lambda_m\)-induced contraction map.
\end{proof}

% ─────────────────────────────────────────────────────────────────────────────
x\section{Implementation Architecture}
% ─────────────────────────────────────────────────────────────────────────────

The thread produced a strong architectural insight: \(\Om\) should not be implemented as a monolithic class alone. Instead, the cleanest structure is a two-object design:
\begin{enumerate}
    \item \texttt{CoreMultiplicityOperator},
    \item \texttt{LawfulRestriction}.
\end{enumerate}

The first object owns the bundle, the global blocks, and the lawfulness projector. The second represents a certified restriction carrying sector metadata, admissibility state, and certification results.

\subsection{Minimal Python Skeleton}

\begin{lstlisting}[style=pythonstyle,caption={Minimal Open Core base class}]
from __future__ import annotations
from dataclasses import dataclass, field
from typing import Any, Dict, Optional
from abc import ABC, abstractmethod

@dataclass
class LawfulRestriction:
    sector: str
    operator: Any
    lawful: bool = True
    metadata: Dict[str, Any] = field(default_factory=dict)
    certificate: Dict[str, Any] = field(default_factory=dict)

    def certify(self, certifier=None) -> Dict[str, Any]:
        if certifier is None:
            return self.certificate
        self.certificate = certifier(self.operator, self.metadata)
        return self.certificate


class OpenCoreSimulator(ABC):
    @abstractmethod
    def build_generators(self, **params) -> Any:
        raise NotImplementedError

    @abstractmethod
    def build_stabilizer(self, **params) -> Any:
        raise NotImplementedError

    @abstractmethod
    def project(self, sector: str, **params) -> LawfulRestriction:
        raise NotImplementedError
\end{lstlisting}

\begin{lstlisting}[style=pythonstyle,caption={Core Multiplicity Operator skeleton}]
from __future__ import annotations
from typing import Any, Dict, Literal, Optional
import numpy as np

Sector = Literal[
    "frame-flat",
    "modular",
    "stochastic",
    "ace-projected",
    "l-functorial",
]

class CoreMultiplicityOperator:
    def __init__(self, sigma: float = 1.0, alpha: float = 1.5,
                 h: Optional[np.ndarray] = None):
        self.sigma = sigma
        self.alpha = alpha
        self.h = h if h is not None else np.ones(1)

        self.prime_id = PrimeIdentity()
        self.langlands = LanglandsRegistrar()
        self.lambdam = LambdaMProtocol()
        self.topos = SovereignTopos()
        self.commission = ToposCommission()

        self.PCSL = self._build_lawfulness_projector()

    def _build_lawfulness_projector(self):
        return self.commission.build_lawfulness_projector(
            self.prime_id, self.lambdam
        )

    def build_core_tensor(self):
        D_sigma = self.prime_id.diag_operator(self.sigma)
        K = self.prime_id.windowed_offdiag(self.alpha, self.h)
        C = self.lambdam.time_sieve_multiplier()
        Xi = self.lambdam.internal_block()
        return D_sigma + K + C + Xi

    def restrict(self, sector: Sector, **params):
        omega = self.build_core_tensor()

        if sector == "frame-flat":
            return LawfulRestriction(sector, omega, metadata=params)

        if sector == "modular":
            op = self.lambdam.moonshine_fixed_point(omega, tau=params.get("tau"))
            return LawfulRestriction(sector, op, metadata=params)

        if sector == "stochastic":
            lam = params.get("dissipation_lambda", 0.0)
            op = self.lambdam.stochastic_noise_kernel(omega, lambda_=lam)
            return LawfulRestriction(sector, op, metadata=params)

        if sector == "ace-projected":
            op = self.lambdam.ace_projected_generator(omega)
            return LawfulRestriction(sector, op, metadata=params)

        if sector == "l-functorial":
            op = self.langlands.prism_from_operator(omega, **params)
            return LawfulRestriction(sector, op, metadata=params)

        raise ValueError(f"Unknown sector: {sector}")

    def certify(self, restriction):
        return self.commission.certify_operator(restriction.operator, self.PCSL)
\end{lstlisting}

\subsection{Why This Design Is Strong}

\begin{itemize}
    \item It preserves existing modules.
    \item It makes lawfulness reusable rather than reimplemented.
    \item It separates mathematical construction from sectoral interpretation.
    \item It is well aligned with future Lean 4 formalization.
\end{itemize}

% ─────────────────────────────────────────────────────────────────────────────
\section{Validation Strategy}
% ─────────────────────────────────────────────────────────────────────────────

The validation path derived in the thread is intentionally staged.

\begin{enumerate}
    \item \textbf{Implement the synthesis layer only.} Do not rewrite existing engines.
    \item \textbf{Check the commutation invariant first.} Treat \([\Om,\PCSL]=0\) as the primary admissibility test.
    \item \textbf{Validate frame-flat restriction first.} This is the cleanest specialization.
    \item \textbf{Then validate modular and stochastic sectors.}
    \item \textbf{Emit machine-readable certificates.} These should correspond to C1--C4.
\end{enumerate}

\begin{remark}
The fastest route to credible validation is not proving everything at once. It is establishing one canonical operator, one projector, one clean restriction, and one successful certificate pathway.
\end{remark}

% ─────────────────────────────────────────────────────────────────────────────
\section{Recommendations}
% ─────────────────────────────────────────────────────────────────────────────

\begin{enumerate}[label=\textbf{R\arabic*.}]
    \item Commit the Multiplicity Open Core specification as a standalone conceptual document.
    \item Implement \texttt{open\_core.py} and \texttt{omega\_operator.py} as separate layers.
    \item Promote lawfulness checks into explicit interfaces instead of keeping them hidden in helper logic.
    \item Add sector-level theorem obligations to every restriction object.
    \item Use the appendix proofs below as the seed for future Lean 4 formalization.
\end{enumerate}

% ─────────────────────────────────────────────────────────────────────────────
\section{Conclusion}
% ─────────────────────────────────────────────────────────────────────────────

The developments of this thread point toward a disciplined unification program. The correct unification is not a premature claim that every sector already is the same operator. Instead, the correct step is to distinguish:
\begin{itemize}
    \item autonomous sectoral engines,
    \item a reusable open-core simulator pattern,
    \item a canonical operator layer,
    \item lawful restrictions,
    \item certification as a machine-checkable witness.
\end{itemize}

This yields a framework that is philosophically coherent, mathematically extensible, software-architecturally clean, and compatible with future formal verification.

% ─────────────────────────────────────────────────────────────────────────────
% ── Appendix counter override (applied at \appendix) ─────────────────────────

 \appendix
 \renewcommand{\thesection}{\Alph{section}}
\renewcommand{\theequation}{\Alph{section}.\arabic{equation}}
   \renewcommand{\thetheorem}{\Alph{section}.\arabic{theorem}}
   \renewcommand{\thelemma}{\Alph{section}.\arabic{lemma}}
   \renewcommand{\theproposition}{\Alph{section}.\arabic{proposition}}
   \renewcommand{\thecorollary}{\Alph{section}.\arabic{corollary}}
   \setcounter{equation}{0}
   \setcounter{theorem}{0}

% 
─────────────────────────────────────────────────────────────────────────────
\section{Proof of the Commutation Invariant}
% ─────────────────────────────────────────────────────────────────────────────

\begin{lemma}
For every prime fibre projection \(\Pi_p\), one has
\[
[D_\sigma,\Pi_p]=0.
\]
\end{lemma}

\begin{proof}
The operator \(D_\sigma\) acts diagonally on each prime fibre \(\mathcal{H}_p\), so projection onto that fibre commutes with \(D_\sigma\).
\end{proof}

\begin{lemma}
If \(\widehat{h}\in L^1(\R)\), then \(K_{\alpha,h}\) is bounded and
\[
\norm{K_{\alpha,h}} \le \norm{\widehat{h}}_{L^1}.
\]
\end{lemma}

\begin{proof}
The block is controlled by its integral kernel representation. The \(L^1\)-bound on \(\widehat{h}\) yields boundedness of the associated convolution-type operator.
\end{proof}

\begin{lemma}
If \(C\) acts only on the \(L^2(\R)\)-factor, then
\[
[C,\Pi_p]=0
\]
for every \(p\in\Pset\).
\end{lemma}

\begin{proof}
The operators act on distinct tensor factors.
\end{proof}

\begin{lemma}
If \(P_{\mathrm{Fix}(\Xi)}\) is the spectral projection onto the fixed-point subspace of \(\Xi\), then
\[
[\Xi, P_{\mathrm{Fix}(\Xi)}]=0.
\]
\end{lemma}

\begin{proof}
Spectral projections commute with the operator from which they arise.
\end{proof}

\begin{proof}[Proof of the Commutation Invariant]
Combine the preceding lemmas and note that \(\PCSL = P_\Pi \wedge P_{\mathrm{Fix}(\Xi)}\). Since each block of \(\Om\) commutes with the relevant projector factor, \([\Om,\PCSL]=0\).
\end{proof}

% ─────────────────────────────────────────────────────────────────────────────
\section{Operator Norm Bounds}
% ─────────────────────────────────────────────────────────────────────────────

\begin{proposition}
\[
\norm{D_\sigma}=|\sigma|.
\]
\end{proposition}

\begin{proof}
\(D_\sigma\) is scalar on each fibre.
\end{proof}

\begin{proposition}
If \(C=\Lambda_m\circ M_\varphi\), then
\[
\norm{C}\le \norm{m}_{L^\infty}\norm{\varphi}_{L^\infty}.
\]
\end{proposition}

\begin{proof}
This follows from submultiplicativity of the operator norm.
\end{proof}

\begin{proposition}
If \(\Xi\) is self-adjoint, then
\[
\norm{\Xi}=r(\Xi).
\]
\end{proposition}

\begin{proof}
For self-adjoint operators, norm equals spectral radius.
\end{proof}

\begin{proof}[Proof of the Global Norm Bound]
Use the triangle inequality on \(\Om = D_\sigma + K_{\alpha,h} + C + \Xi\) and insert the preceding block bounds.
\end{proof}

% ─────────────────────────────────────────────────────────────────────────────
\section{Banach Fixed-Point Certificate}
% ─────────────────────────────────────────────────────────────────────────────

\begin{proof}[Detailed proof of the \(\Lambda_m\) contraction statement]
Let \(\Phi(\psi)=\psi-\lambda_t L\psi\). Then
\[
\norm{\Phi(\psi)-\Phi(\phi)}
=
\norm{(I-\lambda_t L)(\psi-\phi)}
\le
\norm{I-\lambda_t L}\,\norm{\psi-\phi}.
\]
If
\[
\gamma := \sup_t \norm{I-\lambda_t L} < 1,
\]
then \(\Phi\) is a strict contraction. By the Banach fixed-point theorem, \(\Phi\) has a unique fixed point \(\psi^\star\), and iteration converges geometrically:
\[
\norm{\psi_n-\psi^\star}\le \gamma^n \norm{\psi_0-\psi^\star}.
\]
\end{proof}

% ─────────────────────────────────────────────────────────────────────────────
\section{Sector Admissibility}
% ─────────────────────────────────────────────────────────────────────────────

\begin{theorem}
For each admissible sector \(s\), the restriction \(\Om_s=R_s(\Om)\) satisfies
\begin{enumerate}[label=(\roman*)]
    \item \( [\Om_s,\PCSL]=0 \) on the restricted space,
    \item \( \norm{\Om_s}\le \norm{\Om} \),
    \item \(R_{s'}\circ R_s = R_{s'\wedge s}\) whenever the meet exists.
\end{enumerate}
\end{theorem}

\begin{proof}
Restriction is compression to a \(\PCSL\)-invariant subspace, hence preserves commutation and does not increase operator norm. Functoriality follows from composition of inclusions and projections.
\end{proof}

\begin{corollary}
In the frame-flat sector,
\[
\Om_{\mathrm{ff}} = A+B+E,
\]
recovering the Meta-Relativity universal spectral operator in the curvature-free regime.
\end{corollary}

% ─────────────────────────────────────────────────────────────────────────────
\section{Certification Conditions}
% ─────────────────────────────────────────────────────────────────────────────

The four proof obligations attached to a certified restriction are:

\begin{align*}
\textbf{C1: } &\norm{[\Om_s,P]} \le \varepsilon,\\
\textbf{C2: } &\gamma_t \le 1,\\
\textbf{C3: } &R_t \in [\rho_{\min},\rho_{\max}],\\
\textbf{C4: } &\mathcal{M}_t = H(\mathcal{M}_{t-1}\Vert \Om_{s,t}).
\end{align*}

\begin{remark}
These conditions provide a natural interface between runtime validation and proof-assistant formalization.
\end{remark}

% ─────────────────────────────────────────────────────────────────────────────
\section{Bibliography Scaffold}
% ─────────────────────────────────────────────────────────────────────────────

The following works are the most structurally relevant external references for the formulations developed here:

\begin{itemize}
    \item Reed and Simon, \emph{Methods of Modern Mathematical Physics}, Vols. I, II, IV.
    \item Kato, \emph{Perturbation Theory for Linear Operators}.
    \item Kadison and Ringrose, \emph{Fundamentals of the Theory of Operator Algebras}.
    \item Apostol, \emph{Introduction to Analytic Number Theory}.
    \item Connes, \emph{Noncommutative Geometry}; Connes on the trace formula and zeta zeros.
    \item Berry and Keating on spectral interpretations of the Riemann zeros.
    \item Banach, \emph{Théorie des opérations linéaires}.
    \item Langlands, on automorphic forms and functoriality.
    \item Conway--Norton and Borcherds on Monstrous Moonshine.
    \item Mac Lane, Johnstone, Lawvere--Tierney, and Lurie on category and topos theory.
    \item Takens on delay embeddings.
    \item Penrose and Diósi on gravity-related state reduction.
    \item de Moura and Ullrich on Lean 4.
\end{itemize}

\vspace{1em}


\nocite{*}
\bibliographystyle{unsrt}  
\bibliography{references}
\end{document}